\documentclass[11pt,a4paper]{article}
\usepackage[margin=2.2cm]{geometry}
\usepackage[T1]{fontenc}
\usepackage{amsmath,amssymb,xcolor,enumitem,booktabs,graphicx,tabularx,array}
\usepackage[normalem]{ulem}
\setlength{\parskip}{0.6em}\setlength{\parindent}{0pt}
\definecolor{newc}{HTML}{1f6f43}
\newcommand{\new}[1]{\textcolor{newc}{#1}}
\newcommand{\old}[1]{\textcolor{gray}{\sout{#1}}}
\newcommand{\note}[1]{{\small\itshape [#1]}}
\newcommand{\lab}[1]{\textsf{\footnotesize\detokenize{#1}}}
\renewcommand{\ref}[1]{\lab{#1}}
\renewcommand{\eqref}[1]{(\lab{#1})}
\newcommand{\nref}[1]{\lab{#1}}
\newcommand{\citep}[1]{[\lab{#1}]}
\newcommand{\citet}[1]{\lab{#1}}
\newcommand{\pitem}[2]{\subsection*{#1. #2}}
\title{Proposed change, Problem Formulation, version 2\\ (open points after your answers)}
\date{25 November 2026}
\begin{document}\sloppy
\maketitle

\section*{How to read this}
Version 1 (\texttt{review\_problem\_formulation.pdf}) holds P1 to P16. This
version answers your two questions and adds the points of the second list
you pasted. Black text stays, \new{green} is new and \old{struck} text goes.
\note{Notes in brackets} say why. Answer by number, for example ``N3 yes,
N5 keep''.

\textbf{Approved, to apply as in version 1:} P1, P2, P3, P4 option a, P5,
P6, P8, P9, P10, P11, P12, P13, P14. They go in under one tag, [A133]
(another session used [A132] yesterday for the RQ4 sentence in the scope).

\textbf{Still open from version 1:} P7 (answered below, Q1), P15 and P16.
P15 is smaller now. The RQ4 sentence is already cut under [A132], so only
the deadline and re-tasking sentences remain.

\textbf{Already in the thesis.} Two pointer cuts of P1 (after
well-formedness and after Stern et al.) went in before you stopped me. Keep
them, or I undo them until everything goes in together.

% ================================================================
\section*{Your questions}

\pitem{Q1}{P7. Should some of the forward pointers move to the Method?}
No. What the pointers lead to is already in the Method. The two forms of
the distance label moved there in [A124]. The crowding weight is in the
reward table of [A127]. The masked moves are explained in ``How an agent
chooses its move''. So nothing needs to move, only the pointers go.

Since a few forward pointers are allowed, here is every one in the chapter,
with what I propose.

\begin{tabular}{@{}p{0.42\linewidth}p{0.1\linewidth}p{0.4\linewidth}@{}}
\toprule
Pointer & Proposal & Why \\
\midrule
distance label in two forms (\lab{sec:method:model}) & cut & sentence and pointer go, Method has it \\
crowding weight in the reward (\lab{sec:method:rl}) & cut & paragraph goes, reward table has it \\
masked moves (\lab{sec:method:model}) & cut & pointer only, the fact stays \\
\lab{tab:posg}, $R_i$ left open (\lab{sec:method:rl}) & keep & marks something left open on purpose \\
\lab{tab:posg}, $\gamma$ a training value (\lab{sec:method:training}) & keep & same \\
``It is given in Section'' reward (\lab{sec:method:rl}) & keep & same, the one place the text says so \\
scoring opening, ``the reward of Section'' (\lab{sec:method:rl}) & cut & pointer only, the reward is named above \\
the $\Delta$ sweep table & cut & done by P11 \\
\bottomrule
\end{tabular}

That leaves three forward pointers, all to the reward and the discount,
which the chapter leaves open by design. The pointers into the appendices
(stall length, reproduction of Ma et al., crowding comparison, section
field of view) stay. They point to evidence, not to the next step of the
argument.

\pitem{Q2}{P10. Does any of the learning problem need to move to the Method?}
No. What P10 cuts explains a standard game (what a policy is, what a
transition rule is, a probability example), and the chapter needs none of it
again. The one new fact, that the state holds the priority order, stays.
The sequential draw that P10 cuts is told in Who decides and, for the
mask, in the Method's ``Choosing together''. The sentence P10 adds (the
learned policy is the routing rule) belongs here, because it ties the
controller of Who decides to the game.

% ================================================================
\section*{From the second list}
Points of that list that version 1 already covers: 5.3, 5.4 to 5.8, 5.13,
items 1 to 6 and 9, A (P14), B (P10), D (P2) and E (P5). Below are the ones
that are new, or where the list differs from version 1.

\pitem{N1}{``Storage epochs'' (Fix)}
In \emph{How the two layers close the loop}:

Only every $\Delta$ timesteps, at the storage \old{epochs}\new{updates}, does
new stock arrive and the storage rule place it.

\note{``Epoch'' is the old word. [A123] renamed it to update everywhere
else.}

\pitem{N2}{``Queue split'' in the measures table (Fix)}
In \lab{tab:measures}, the source column of the split:

\old{queue split}\new{service-time split}, this thesis

\note{One name per idea. The text calls it the service-time split.}

\pitem{N3}{The Introduction on demand storage (Fix)}
Demand storage puts the SKUs ordered most lately on the free shelves
\old{nearest the workstations}\new{with the shortest trips}.

\note{The code and the Method rank a shelf by the trip from a dock, through
the shelf, to a workstation, with the pick leg weighted by how often it is
walked [A122]. ``Nearest the workstations'' leaves the dock out. This is in
the Introduction, which another session worked on yesterday ([A130],
[A132]). It is one phrase and does not touch their changes.}

\pitem{N4}{The storage update example in one sentence (Choice)}
\old{Take the stock of Table~\ref{tab:storagematrix} just before an update,
with nothing on its way, and say three pallets of tea arrive. For the
example a pallet holds 3 units. Shelf $A$ is full, and $B$ and $C$ each have
room for one more pallet. The rule puts the first pallet on $B$ and the
second on $C$, and both shelves stay within their room. Nothing was on its
way, so no earlier pallet had to keep its shelf. The third pallet fits
nowhere and is left unplaced, so the rule places 6 of the 9 units that
arrived. Figure~\ref{fig:storageupdate} shows the update.}
\new{Figure~\ref{fig:storageupdate} shows an update on the stock of
Table~\ref{tab:storagematrix}, when three pallets of 3 tea arrive. Shelf $A$
is full, so two pallets go to $B$ and $C$, and the third fits nowhere and is
lost.}

\note{The figure and its caption already show each step. The paragraph
named all three conditions one by one. Two of them (nothing on its way,
only arriving pallets) do nothing in this example. Recommended.}

\pitem{N5}{The subtraction $\zeta_1 = 14-5 = 9$ (Choice)}
The second list says drop it. Version 1 kept it, since your style asks for
a worked number after a formula and it is one line. My recommendation is
still to keep it. If you want it gone:

The service time of task $\tau_j$ is the time from its release until its
delivery, as in \citet{ma2017lifelong_mapf}. Formally, $\zeta_j =
t_j^{\mathrm{finish}} - t_j^{\mathrm{released}}$. \old{For $\tau_1$, the
service time is $\zeta_1=14-5=9$.}

\pitem{N6}{The $2/5$ crowding example (Choice)}
You approved P8, which keeps it as one sentence. The second list says
drop it entirely, since \lab{fig:crowding} gives an example right after.
The one sentence ties the feature to the figure the section already uses,
so I recommend P8 as approved. If you want it gone:

\old{For agent $a_1$ in Figure~\ref{fig:observation}, two of the five other
vertices in view hold an agent, so $\delta_1(t) = 0.4$.}

\pitem{N7}{The crowding feature under Who decides (Choice)}
Every controller reads the crowding feature, but it is defined under
\emph{What a decentralised agent sees}. The second list asks whether it
should move to Who decides.

\textbf{Recommended: keep it where it is.} The feature is a share of the
field of view, so it needs the field of view first, which is defined in
the same section. Moving it would pull the field of view into Who decides
too. The section already says that every controller reads it. The only
change is a clearer heading for the subsection:

\old{How crowded the neighbourhood is}\new{How crowded the neighbourhood is,
for every controller}

\note{Option b: move the subsection and the field of view to the end of Who
decides, and keep only the target, messages and the observation equation
in ``What a decentralised agent sees''. More moving for little gain.}

% ================================================================
\section*{Still open from version 1}

\pitem{P15}{Scope, the boundary only}
Version 1 as before, without the RQ4 part, which [A132] already cut. The
second list (item 9) proposes the same.

\old{Adding deadlines would put a deadline in every task, and either forbid
a late task or count late tasks next to the result.}

\old{Adding it would let the assignment rule change the agent of a task
whose agent has not yet reached the shelf. Either change alters what counts
as a solution, so it would change this chapter and not the method.}
\new{Either change would alter what counts as a solution.}

\pitem{P16}{Old text in comments}
About 200 lines of commented-out old text in the chapter, not in the PDF,
kept in git. Delete them?

% ================================================================
\section*{When approved}
Everything goes in under [A133] with one GAPS.md entry. P4a changes the
row of $\mathcal{F}$ in both notation tables. P14 moves
\lab{eq:rqformal} and \lab{eq:relgain} to the Experiments. Their labels stay,
so the references in the Experiments and Appendices B, F and K still work
after \texttt{tools/sync\_xrefs.py --build}. Every tag goes at the end of a
line. The chapter is rebuilt and read back with pdftotext.

\end{document}
